%! TEX root = ../bcd-19-20.tex

\chapter{Extreme libere}

\section{Polinomul Taylor în 2 variabile}

Fie $ f : D \seq \RR^2 \to \RR $ o funcție de două variabile,
de clasă $ \kal{C}^\infty $. \emph{Polinomul Taylor} al lui $ f $
în jurul punctului $A(x_0, y_0) $ se scrie:
\begin{align*}
    T(x, y) &= f(x_0, y_0) + \frac{1}{1!} \Big[ (x - x_0) \frac{\partial f}{\partial x}(A) + %
    (y - y_0) \frac{\partial f}{\partial y}(A) \Big] + \\
            &+ \frac{1}{2!} \Big[ (x - x_0)^2 \frac{\partial^2 f}{\partial x^2}(A) + %
                2(x - x_0)(y - y_0) \frac{\partial^2 f}{\partial x \partial y}(A) + %
            (y - y_0)^2 \frac{\partial^2 f}{\partial y^2}(A) \Big] + \\
            &+ \frac{1}{3!} \Big[ (x - x_0)^3 \frac{\partial^3 f}{\partial x^3}(A) + %
                3(x - x_0)^2(y - y_0) \frac{\partial^3 f}{\partial^2 x \partial y}(A) + \\
            &+ 3(x - x_0)(y - y_0)^2 \frac{\partial^3 f}{\partial x \partial^2 y}(A) + %
            (y - y_0)^3 \frac{\partial^3 f}{\partial y^3}(A) \Big] \\
            &+ \dots
\end{align*}

Gradul polinomului Taylor este dat de ordinul ultimei derivate parțiale
care se calculează. Astfel, avem, în particular:
\begin{align*}
  T_1(x, y) &= f(x_0, y_0) + \frac{1}{1!} \Big[ %
      (x - x_0) \frac{\partial f}{\partial x}(A) + %
  (y - y_0) \frac{\partial f}{\partial y}(A) \Big] \\
  T_2(x, y) &= T_1(x, y) + \frac{1}{2!} \Big[ %
      (x - x_0)^2 \frac{\partial^2 f}{\partial x^2}(A) + %
      2(x - x_0)(y - y_0) \frac{\partial^2 f}{\partial x \partial y}(A) + %
  (y - y_0)^2 \frac{\partial^2 f}{\partial y^2}(A) \Big].
\end{align*}

%%%%%%%%%%%%%%%%%%%%%%%%%%%%%%%%%%%%%%%%%%%%%%%%%%%%%%%%%%%%%%%%%%%%%%

\section{Extreme libere}

Fie $ f : D \seq \RR^2 \to \RR $ o funcție de două variabile. Pentru a
găsi punctele de extrem ale funcției, parcurgem următoarele etape:
\begin{enumerate}[(1)]
\item rezolvăm sistemul de ecuații dat de anularea derivatelor parțiale:
  \[
    \begin{cases}
      \dfrac{\partial f}{\partial x} &= 0 \\
      & \\
      \dfrac{\partial f}{\partial y} &= 0
    \end{cases}
  \]
  Soluțiile acestui sistem se numesc \emph{puncte critice} ale funcției.
\item scriem \emph{matricea hessiană} a funcției, adică:
  \[
    H_f =
    \begin{pmatrix}
      \dfrac{\partial^2 f}{\partial x^2} & \dfrac{\partial^2 f}{\partial x \partial y} \\
      & \\
      \dfrac{\partial^2 f}{\partial y \partial x} & \dfrac{\partial^2 f}{\partial y^2}
    \end{pmatrix}
  \]
\item Evaluăm matricea în fiecare punct critic. Astfel, dacă $ A(x_0, y_0) $ este
  punct critic, calculăm $ H_f(x_0, y_0) $;
\item Determinăm \emph{valorile proprii} $ \lambda_{1,2} $ ale matricei hessiene,
  adică rădăcinile polinomului caracteristic $ P_M(X) = \det (M - XI_2) $,
  unde $ M = H_f(x_0, y_0)$.
  \begin{itemize}
  \item Dacă $ \lambda_1, \lambda_2 > 0 $, atunci punctul $ A(x_0, y_0) $
    este \emph{punct de minim local};
  \item Dacă $ \lambda_1, \lambda_2 < 0 $, atunci punctul $ A(x_0, y_0) $
    este \emph{punct de maxim local};
  \item Dacă $ \lambda_1 \cdot \lambda_2 < 0 $, atunci punctul $ A(x_0, y_0) $
    \emph{nu este de extrem};
  \item Dacă $ \lambda_1 \cdot \lambda_2 = 0 $, problema necesită un studiu
    separat, acesta fiind un caz mai dificil și pe care îl omitem. Formal, trebuie
    studiat semnul diferenței $ f(x, y) - f(x_0, y_0) $, pentru fiecare punct
    critic $ A(x_0, y_0) $, iar pentru aceasta se folosește polinomul Taylor.
  \end{itemize}
\item Se repetă procedura pentru fiecare punct critic.
\end{enumerate}

Procedura de mai sus se aplică în cazul așa-numitelor \emph{extreme libere},
adică atunci cînd domeniul este $ D = \RR^2 $ sau un dreptunghi dat de un
produs cartezian de intervale.

Pentru cazul cînd domeniul este dat de (in)ecuații, procedura este diferită
și o vom exemplifica mai jos.

%%%%%%%%%%%%%%%%%%%%%%%%%%%%%%%%%%%%%%%%%%%%%%%%%%%%%%%%%%%%%%%%%%%%%%

\section{Exerciții rezolvate}

Fie funcția $ f : D \seq \RR^2 \to \RR, f(x, y) = x^3 + y^3 - 6xy $.
\begin{enumerate}[(a)]
\item Pentru $ D = \RR^2 $, determinați punctele de extrem și calculați
  valorile funcției în aceste puncte.
\item Pentru $ D = \{ (x, y) \in \RR^2 \mid x \geq 0, y \geq 0, x + y \leq 5 \} $,
  determinați punctele de extrem și valorile extreme.
\end{enumerate}

\emph{Soluție:}

(a) Punctele critice se determină din anularea derivatelor parțiale. Avem:
\[
  \frac{\partial f}{\partial x} = 3x^2 - 6y, \quad %
  \frac{\partial f}{\partial y} = 3y^2 - 6x.
\]
Se obțin $ A(0, 0) $ și $ B(2, 2) $.

Apoi matricea hessiană:
\[
  H_f = \begin{pmatrix}
    6x & -6 \\
    -6 & 6y
  \end{pmatrix}
\]

Pentru $ A(0, 0) $, hessiana este:
\[
  H_f(A) = \begin{pmatrix}
    0 & -6 \\
    -6 & 0
  \end{pmatrix} = M.
\]

Polinomul caracteristic $ P_M(X) = \det (M - XI_2) = X^2 - 36 \Rightarrow \lambda_{1,2} = \pm 6 $.
Rezultă că punctul $ A $ nu este de extrem.

Pentru $ B(2, 2) $, hessiana este:
\[
  H_f(B) = \begin{pmatrix}
    12 & -6 \\
    -6 & 12
  \end{pmatrix} = M.
\]

Polinomul caracteristic $ P_M(X) = (12 - X)^2 - 36 \Rightarrow \lambda_1 = 6, \lambda_2 = 18 $.
Cum ambele valori proprii sînt pozitive, punctul $ B $ este punct de minim local.

Valoarea minimă a funcției este $ f(2, 2) = -8 $.

\vspace{1cm}

(b) Problema se studiază pe 2 cazuri: pe interiorul lui $ D $ și pe frontiera lui $ D $.

Pentru frontieră, avem $ x \geq 0, y \geq 0 $ și $ x + y = 5 $. Atunci putem
studia funcția $ g(x) = f(x, 5-x) = 21x^2 - 105x + 125 $, care are un punct
de minim în vîrful acestei parabole. Se obține $ f(2,5; 2,5) = -6,25 $.
De asemenea, mai avem și valorile $ f(0, 5) = 125 $, respectiv $ f(5, 0) = 125 $.

Apoi mai avem de studiat cazurile $ y = 0, x \in [0, 5] $, deci funcția
$ f(x, 0) = h(x) $, precum și $ x = 0, y \in [0, 5] $, deci funcția
$ f(0, y) = j(y) $.

Pentru interiorul domeniului, avem:
\[
  \dr{Int}(D) = \{ (x, y) \in \RR^2 \mid x > 0, y > 0, x + y < 5.
\]

Putem folosi calculele din cazul anterior și constatăm că $ B(2, 2) \in \dr{Int}(D) $,
deci avem o valoare de minim pe interior, cu $ f(2, 2) = -8 $.

Așadar, răspunsul final este $ \max f = 125 $ și $ \min f = -8 $.

%%%%%%%%%%%%%%%%%%%%%%%%%%%%%%%%%%%%%%%%%%%%%%%%%%%%%%%%%%%%%%%%%%%%%%

\section{Exerciții propuse}

1. Determinați punctele de extrem și valorile extreme pentru funcțiile:
\begin{enumerate}[(a)]
\item $ f : D \seq \RR^2 \to \RR, f(x, y) = x^2 + xy + y^2 - 4 \ln x - 10\ln y + 3 $;
\item $ f(x, y) = 3x^2 - x^3 - 15x - 36y + 9 $, definită (i) pe $ \RR^2 $
  și (ii) pe $ [-4,4] \times [-3, 3] $;
\item $ f(x, y) = 4xy - x^4 - y^4 $, definită (i) pe $ \RR^2 $ și (ii) pe
  $ [-1, 2] \times [0, 2] $;
\item $ f(x, y) = x^3 + 3x^2y - 15x - 12y $ definită (ii) pe $ \RR^2 $ și
  (ii) pe $ D = \{ x \geq 0, y \geq 0, 3y + x \leq 3 \} $;
\item $ f(x, y) = xy(1 - x - y) $, definită pe $ [0, 1] \times [0, 1] $;
\item $ f(x, y) = x^3 + 8y^3 - 2xy $, definită pe $ D = \{ x, y \geq 0, y + 2x \leq 2 \} $;
\item $ f(x, y) = x^4 + y^3 - 4x^3 - 3y^2 + 3y $, definită pe
  $ D = \{ x^2 + y^2 \leq 4 \} $.
\end{enumerate}
 
În toate exercițiile de mai sus este recomandabil să reprezentați
grafic domeniile de definiție.

\vspace{1cm}

2. Fie funcția $ f(x,y) = x \ln(x^2 + y^2) $. Calculați polinoamele Taylor
de gradul 1 și respectiv 2, $ T_1, T_2 $ în jurul punctului $ (1, 0) $.

\vspace{1cm}

3. Aceeași cerință pentru funcția $ f(x, y) = 3xe^{x^2 + y^2} $ și pentru
funcția $ g(x, y) = \arctan \dfrac{x}{y} $.


